% !TeX program = pdflatex
% A proof of Kurkov's conjecture for A000670
% Build: pdflatex note.tex (run twice for cross-references).
\documentclass[11pt]{article}

\usepackage[T1]{fontenc}
\usepackage{lmodern}
\usepackage[margin=1in]{geometry}
\usepackage{amsmath,amssymb,amsthm}
\usepackage{microtype}
\usepackage[hidelinks]{hyperref}
\usepackage{xurl}

\newtheorem{theorem}{Theorem}[section]
\newtheorem{lemma}[theorem]{Lemma}
\newtheorem{corollary}[theorem]{Corollary}
\theoremstyle{definition}
\newtheorem{definition}[theorem]{Definition}
\theoremstyle{remark}
\newtheorem{remark}[theorem]{Remark}

\DeclareMathOperator{\wt}{wt}
\newcommand{\bit}{\mathbf{b}}

\title{A proof of Kurkov's conjecture for A000670}
\author{John Erlbacher\thanks{Independent Researcher. Email:
\texttt{erlbacher.research@gmail.com}. ORCID: 0009-0003-6851-4139.}}
\date{July 2026}

\begin{document}
\maketitle

\begin{abstract}
The Fubini numbers $a(n)$ (OEIS \texttt{A000670}), which count the ordered set
partitions of an $n$-element set, satisfy a conjecture of Mikhail Kurkov (2018)
recorded on the entry: for $n>0$,
\[
  a(n)=\sum_{k=0}^{2^{\,n-1}-1} A284005(k),\qquad a(0)=1,
\]
where $A284005$ (OEIS \texttt{A284005}) is defined by $A284005(0)=1$ and
$A284005(m)=(1+\wt(m))\,A284005(\lfloor m/2\rfloor)$, with $\wt$ the binary digit
sum. We give a short, elementary, and self-contained proof. The key is a refined
statement: for each of the $2^{\,n-1}$ subsets $S\subseteq[n]$ containing $1$,
the number of ordered set partitions of $[n]$ whose set of block minima equals
$S$ is exactly $A284005(k)$ for the index $k$ that encodes $S$. Summing over $S$
proves the conjecture. The refinement is verified exhaustively for $n\le 8$
($598{,}444$ ordered set partitions across $255$ block-minima classes, $0$
mismatches).
\end{abstract}

\section{Introduction}\label{sec:intro}

Let $a(n)$ denote the number of \emph{ordered set partitions} of $[n]=\{1,\dots,n\}$:
partitions of $[n]$ into nonempty blocks together with a linear order on the
blocks. These are the Fubini (or ordered Bell) numbers, OEIS entry
\texttt{A000670} \cite{A000670}, with $a(0)=1$ and initial values
$1,1,3,13,75,541,4683,\dots$

The entry \texttt{A000670} carries the following formula, contributed by
Mikhail Kurkov on July~8, 2018, and still labelled a conjecture:
\begin{equation}\label{eq:conj}
  a(n)=\sum_{k=0}^{2^{\,n-1}-1} A284005(k)\qquad(n>0),\qquad a(0)=1,
\end{equation}
where $A284005$ is the OEIS entry \texttt{A284005} \cite{A284005}, defined by
\begin{equation}\label{eq:a284005}
  A284005(0)=1,\qquad A284005(m)=(1+\wt(m))\,A284005(\lfloor m/2\rfloor),
\end{equation}
and $\wt(m)$ is the number of $1$'s in the binary expansion of $m$. The purpose
of this note is to prove \eqref{eq:conj}.

The proof is elementary and combinatorial. We first unwind \eqref{eq:a284005}
into a product over the binary digits of $k$ (Section~\ref{sec:setup}); this
lets us index the summands of \eqref{eq:conj} by subsets of $[n]$ containing
$1$. We then prove a refinement of \eqref{eq:conj} (Section~\ref{sec:proof}):
each summand counts exactly the ordered set partitions with a prescribed set of
block minima. Because every ordered set partition of $[n]$ has a unique set of
block minima, and $1$ is always a block minimum, summing the refinement over all
$2^{\,n-1}$ admissible subsets yields \eqref{eq:conj}.

To our knowledge \eqref{eq:conj} was open before this work: it is stated as a
conjecture on \texttt{A000670}, and a dated literature and web search
(July~2026) turned up no proof.

\section{The product formula and its index set}\label{sec:setup}

Iterating \eqref{eq:a284005} strips the binary expansion of $m$ from the least
significant bit. Writing $m$ in binary as $b_1b_2\cdots b_L$ (most significant
bit first, $b_1=1$), the successive arguments $\lfloor m/2\rfloor,
\lfloor m/4\rfloor,\dots$ are the prefixes $b_1\cdots b_{L-1}, b_1\cdots
b_{L-2},\dots$, and each contributes the factor $1+(\text{its digit sum})$.
Hence
\begin{equation}\label{eq:product}
  A284005(m)=\prod_{i=1}^{L}\bigl(1+w_i\bigr),\qquad
  w_i:=b_1+b_2+\cdots+b_i,
\end{equation}
the same product formula recorded on \texttt{A284005}. A prefix of digit sum
$0$ contributes a factor $1+0=1$; equivalently, \emph{prepending a $0$ to the
binary string does not change the value}. (Here $A284005(0)$ is the empty
product $1$.)

Fix $n>0$. Every index $k$ with $0\le k\le 2^{\,n-1}-1$ can be written, after
left-padding with zeros, as a string of exactly $n-1$ bits
$\bit=b_1b_2\cdots b_{n-1}$ (most significant first). By the padding invariance
just noted, \eqref{eq:product} gives
\begin{equation}\label{eq:kproduct}
  A284005(k)=\prod_{i=1}^{n-1}\bigl(1+w_i\bigr),
  \qquad w_i=b_1+\cdots+b_i,
\end{equation}
for every such $k$, using its $(n-1)$-bit padded expansion. As $k$ ranges over
$\{0,\dots,2^{\,n-1}-1\}$, the string $\bit$ ranges over all of
$\{0,1\}^{\,n-1}$.

\begin{definition}\label{def:M}
For $\bit=b_1\cdots b_{n-1}\in\{0,1\}^{\,n-1}$ put
\[
  M(\bit)=\{1\}\cup\{\,i+1 : b_i=1\,\}\subseteq[n].
\]
\end{definition}

The map $\bit\mapsto M(\bit)$ is a bijection from $\{0,1\}^{\,n-1}$ onto
$\{S\subseteq[n] : 1\in S\}$: bit $b_i$ records whether $i+1\in S$, and $1$ is
always included. There are $2^{\,n-1}$ such subsets, matching the range of $k$.

\section{The refined theorem}\label{sec:proof}

\begin{theorem}\label{thm:refined}
Let $n>0$ and $\bit\in\{0,1\}^{\,n-1}$. The number of ordered set partitions of
$[n]$ whose set of block minima equals $M(\bit)$ is exactly
$\prod_{i=1}^{n-1}(1+w_i)$, where $w_i=b_1+\cdots+b_i$.
\end{theorem}

Before proving Theorem~\ref{thm:refined} we record how it settles the
conjecture.

\begin{corollary}
For every $n>0$, $a(n)=\sum_{k=0}^{2^{\,n-1}-1}A284005(k)$.
\end{corollary}

\begin{proof}
Every ordered set partition of $[n]$ has a well-defined set of block minima,
which necessarily contains $1$ (the block containing $1$ has minimum $1$). Thus
the ordered set partitions of $[n]$ are partitioned according to their
block-minima set $S$, and $S$ ranges over exactly $\{S\subseteq[n]:1\in S\}$.
By Theorem~\ref{thm:refined} and \eqref{eq:kproduct}, the class with minima set
$S=M(\bit)$ has size $A284005(k)$, where $k$ is the integer with padded binary
string $\bit$. Summing over the $2^{\,n-1}$ admissible $S$, equivalently over
$k=0,\dots,2^{\,n-1}-1$, gives $a(n)=\sum_k A284005(k)$. The case $n=0$ is the
boundary value $a(0)=1$.
\end{proof}

\begin{proof}[Proof of Theorem~\ref{thm:refined}]
We build the ordered set partitions of $[n]$ by inserting the elements
$1,2,\dots,n$ one at a time in increasing order, tracking the ordered list of
blocks. We use one observation throughout: \emph{because elements are inserted
in increasing order, an element is the minimum of its block if and only if it
opened that block} --- every element that joins an already existing block is
larger than that block's opener, hence not its minimum.

Start with element $1$, which necessarily opens the single block $(\{1\})$;
after this step there is $m_1=1$ block. Now suppose we have an ordered set
partition of $[i]$ with $m_i$ blocks, and we insert element $i+1$
($1\le i\le n-1$). We want element $i+1$ to be a block minimum precisely when
$b_i=1$ (so that, by Definition~\ref{def:M}, $i+1$ enters the minima set exactly
when $b_i=1$). By the observation above:
\begin{itemize}
\item if $b_i=1$ (element $i+1$ must be a block minimum), it must \emph{open a
  new singleton block}, which can be inserted in any of the $m_i+1$ positions of
  the ordered list of blocks;
\item if $b_i=0$ (element $i+1$ must not be a block minimum), it must \emph{join
  one of the $m_i$ existing blocks}.
\end{itemize}
The number of blocks after step $i$ is $m_{i+1}=1+w_i$: it starts at
$m_1=1=1+w_0$ (with $w_0=0$) and increases by one exactly when $b_i=1$, i.e.\
$m_{i+1}=m_i+b_i$ and $w_i=w_{i-1}+b_i$. Consequently the number of choices at
step $i$ is, in both cases, exactly $1+w_i$:
\[
  b_i=1:\quad m_i+1=(1+w_{i-1})+1=1+w_i\ \ (w_i=w_{i-1}+1);\qquad
  b_i=0:\quad m_i=1+w_{i-1}=1+w_i\ \ (w_i=w_{i-1}).
\]
Therefore the number of insertion histories realising the prescribed
signature $\bit$ is $\prod_{i=1}^{n-1}(1+w_i)$.

It remains to see that these insertion histories are in bijection with the
ordered set partitions of $[n]$ whose block-minima set is $M(\bit)$. Distinct
choices at any step produce distinct ordered set partitions, so distinct
histories give distinct outputs. Conversely, every ordered set partition $P$ of
$[n]$ arises from exactly one history: deleting the element $n$ (and deleting its
block if $n$ was a singleton) recovers an ordered set partition of $[n-1]$
together with the last choice made, and this reconstruction is deterministic;
induction on $n$ gives a unique history for $P$. Under this bijection the
signature of the history recording which insertions opened a new block is exactly
the indicator of the block minima of $P$ among $2,\dots,n$ --- since removing the
maximum of a block never changes that block's minimum --- so the histories with
signature $\bit$ correspond precisely to the ordered set partitions with
block-minima set $M(\bit)$. This proves the count.
\end{proof}

\section{Verification}\label{sec:verify}

The identity \eqref{eq:conj} was checked numerically for $n=1,\dots,20$ by
computing $a(n)$ from the Fubini recurrence and independently from the
right-hand side of \eqref{eq:conj} via \eqref{eq:a284005}; the two agree.

Theorem~\ref{thm:refined} was verified \emph{exhaustively} for $n=1,\dots,8$ by
generating every ordered set partition of $[n]$ (there are
\[
  1+3+13+75+541+4683+47293+545835=598{,}444
\]
of them in total across $n=1,\dots,8$), grouping them by block-minima set, and
comparing the class sizes to $A284005(k)$ over all $255$ minima-set classes
($1+2+\cdots+128$). There were $0$ mismatches. A mutation-testing harness, which
perturbs the $A284005$ recurrence and the block-minima encoding, confirms that
the check is a genuine discriminator: every mutant is rejected.

These computations corroborate but do not replace the proof of
Section~\ref{sec:proof}, which is elementary, finite, and self-contained.

\section*{Acknowledgments}

This note, including all of its mathematical content, was produced by an
artificial-intelligence workflow, and we believe readers and referees should be
able to take this fact into account. The conjecture was selected, attacked, and
proved by language-model agents (Anthropic's Claude family) running in
Demonstrandum, an orchestrated, verification-first multi-agent research
pipeline. The definitions of \texttt{A000670} and \texttt{A284005} were pinned
as verbatim quotations of the OEIS entries before any proof attempt; the proof
was then subjected to an adversarial audit in which a model from a different
family (OpenAI's GPT-5.5, accessed through the Codex CLI with code-execution
access) acted as a hostile referee, prompted to find counterexamples and gaps
rather than to confirm, and wrote its own clean-room verification code. No human
mathematician contributed mathematical content; the human role was limited to
operating the pipeline. Correctness rests on the elementary proof printed above,
intended to be checkable line by line in the ordinary way; the machine checks of
Section~\ref{sec:verify} are corroboration, not a substitute.

\begin{thebibliography}{9}

\bibitem{A000670}
OEIS Foundation Inc.,
Entry \texttt{A000670}, \emph{Fubini numbers} (ordered set partitions), in
\emph{The On-Line Encyclopedia of Integer Sequences},
\url{https://oeis.org/A000670} (accessed July 2026). Conjecture of
M.~Kurkov, July~8, 2018.

\bibitem{A284005}
OEIS Foundation Inc.,
Entry \texttt{A284005}, in \emph{The On-Line Encyclopedia of Integer
Sequences}, \url{https://oeis.org/A284005} (accessed July 2026).

\end{thebibliography}

\end{document}
